% 08-optimization.tex

\section{Model-Aware Optimization}\label{sec:optimization}

Designs drawing on established ML techniques.
Crucible's contribution is the infrastructure: each optimization is a DAG branch (Section~3.4), verified before activation, rollbackable via atomic swap.

\subsection{Token-Level Adaptation}

\textbf{Token merging.}
Pairwise cosine similarity between adjacent representations after layer~$N$; tokens above threshold are merged (averaged).
Attention is $O(n^2)$; $4\times$ fewer tokens $= 16\times$ less attention.
Threshold adapts per input and per layer [Bolya et al.\ 2023].

\textbf{Early exit.}
$\|h_N - h_{N-1}\|$ per token below threshold $\to$ freeze, skip remaining layers.
Tokens bucketed into convergence groups for batch efficiency.

\textbf{Adaptive patching (images).}
Quadtree decomposition by information content.
Low-information regions get fewer, larger patches.

Each adaptation is a BranchNode with the non-adapted path as fallback arm.

\textbf{Per-token precision.}
High-information tokens run in FP16; low-information tokens in INT8 or INT4.
Separate kernels per precision group, dispatched from measured information content per token per layer.

\subsection{Layer-Level Analysis}

Operations grouped by \code{ScopeHash}; each layer analyzed independently via Augur diagnostics (Section~6.4).

\textbf{Attention head classification.}
From recorded attention matrices: positional heads (diagonal band) $\to$ depthwise convolution $O(n \cdot k)$ vs $O(n^2)$; global heads (fixed landmarks) $\to$ gather + broadcast $O(n)$; averaging heads (uniform) $\to$ mean pooling $O(n)$; dead heads (near-zero output) $\to$ removed.
Content-routing heads retained or replaced with hash-based routing.
For a 144-head transformer, typical result: ${\sim}85$ sparse-attention + 15 conv + 20 pool + 10 removed + 14 routing.
Total attention cost drops ${\sim}60\%$.

\textbf{Local learning signals.}
Per-layer loss functions (predictive coding, contrastive, reconstruction) inserted as DAG modifications, providing gradient signal with depth-1 backpropagation paths.

\textbf{Per-layer gradient strategy.}
From gradient SNR, Jacobian rank, gradient norms: high SNR $\to$ standard backprop; moderate SNR $\to$ K-FAC natural gradient [Martens and Grosse 2015]; near-zero SNR $\to$ synthetic gradients; converged $\to$ frozen.
50--70\% of layers skippable late in training.

\textbf{Matrix structure discovery.}
Not every dense matmul needs to be dense.
Crucible observes weight matrices per layer and classifies: full-rank (keep dense); low-rank (effective rank $r \ll d$, replace $W(d{\times}d)$ with $A(d{\times}r) \cdot B(r{\times}d)$, $2\times$ cheaper at $r = d/4$); near-Toeplitz (diagonal structure, replace with depthwise conv + small correction, $10\times$ cheaper); highly sparse ($>$95\% near-zero, cuSPARSE); block-diagonal (independent subspace matmuls, naturally parallelizable across streams).
Each replacement is a DAG branch verified against the original.

\textbf{NaN/Inf early kill.}
Lightweight \code{isfinite} checks at numerically sensitive points (softmax, log, exp, division)---${\sim}1$\,\textmu{}s per check.
The moment a NaN appears, Crucible catches it before propagation through 50 more ops.
Response: rollback to previous iteration parameters, skip bad batch, continue.
Current PyTorch: silent NaN propagation, user notices 20 minutes later.

\subsection{Architecture Mutation}

The DAG IS the architecture.
Every modification is a BranchNode: old architecture (arm~A) vs new (arm~B), verified on validation data, committed or discarded via atomic swap.

\textbf{Layer growing.}
Loss plateau + capacity analysis $\to$ insert layer at highest-gradient position.
Initialize via identity/distillation/random.

\textbf{Layer pruning.}
CKA $> 0.95$ between adjacent layers $\to$ skip redundant layer.

\textbf{Width mutation.}
Effective rank $\ll$ hidden dimension $\to$ reduce width via SVD projection + adapters.
A 4096-dim layer with effective rank 600 $\to$ insert $W_\text{down}(4096{\times}600) \cdot W_\text{up}(600{\times}4096)$, ${\sim}3.4\times$ cheaper.

\textbf{Progressive growing.}
Start small (e.g.\ 4 layers, $d{=}512$), grow on plateaus, widen when rank saturates, prune when redundant.
The model's size trajectory is determined by data and loss landscape, not by human guess at step~0.
Typical trajectory: 4~layers at step~0 $\to$ 6 at 10K (loss plateau) $\to$ 8~layers, $d{=}1024$ at 30K (depth needs capacity) $\to$ stable at 100K $\to$ pruned to 9~layers at 150K (dead layer removed).

\textbf{Model composition.}
Two pre-trained models composed by DAG splicing: connect vision encoder DAG to language model DAG with an adapter RegionNode.
Both subgraphs retain \code{content\_hashes}; compiled kernels reused from the KernelCache.
Only the adapter needs compilation.

\textbf{Activation function evolution.}
Per-layer: try SwiGLU, ReLU, GELU as branch arms, measure quality impact.
Typical result: SwiGLU for early layers (expressiveness), ReLU for middle (cheaper, sufficient), GELU for final (refinement).

\subsection{Training Optimization}

The entire training loop is in the DAG and therefore observable and modifiable.

\textbf{Meta-gradients.}
$\partial(\text{val\_loss})/\partial(\text{lr})$ via one additional backward pass.
Learning rate, weight decay, momentum parameters tune themselves by gradient descent on validation loss.
No search.

\textbf{Per-layer LR from curvature.}
Hessian diagonal (periodic Hessian-vector products) gives per-parameter curvature.
Optimal lr proportional to $1/H_{ii}$.

\textbf{Curriculum learning.}
Per-sample loss observable during recording.
Order by difficulty (easy-to-hard, hard-first, random).
The ordering is a DAG-level decision evaluated empirically.

\textbf{Automatic mixed precision.}
Run each op in FP32, BF16, FP8; measure per-op quality impact; select cheapest precision within tolerance.
Per-op, per-training-stage.
Not a static allow-list---discovered from this model on this data.

\textbf{Manifold Mixup.}
Interpolate hidden states between samples at intermediate layer~$K$: $h_\text{mix} = \alpha \cdot h_A + (1 - \alpha) \cdot h_B$.
Forward remainder, compute loss against interpolated label.
Layer~$K$ chosen by linear probe accuracy.
Generates training signal from latent geometry without new data.

\textbf{Representation steering (inference).}
Compute direction vectors in latent space: mean truthful activations minus mean hallucinated activations.
Add $\alpha \times \text{direction}$ at optimal layer during inference.
No weight changes; one vector addition per layer, negligible cost, measurable behavior modification.
Different deployment contexts use different steering vectors.

\textbf{Loss function evolution.}
Auxiliary losses (contrastive at layer~6, regularization against representation collapse) weighted by meta-gradients: $\partial(\text{val\_loss})/\partial(\text{aux\_weight})$ determines whether the auxiliary loss helps.

\textbf{Optimizer evolution.}
Adam, AdaFactor, Lion, and learned update rules as DAG branches.
Crucible evaluates each on validation slices and activates the winner.
The optimizer itself becomes a tunable component.

\subsection{Verification and Activation}

Figure~\ref{fig:optloop} illustrates the uniform protocol.

\begin{figure}[ht]
\centering
\begin{tikzpicture}[
  stp/.style={draw=black!30, fill=white, rounded corners=2pt,
    minimum width=9em, minimum height=2.2em,
    font=\sffamily\scriptsize, align=center,
    blur shadow={shadow blur steps=3, shadow xshift=0.3pt,
                 shadow yshift=-0.3pt, shadow blur radius=1pt}},
  arr/.style={-{Stealth[length=4.5pt]}, thick, black!35},
  lbl/.style={font=\sffamily\tiny, text=black!40},
]
  %% Steps
  \node[stp, fill=blue!5] (diag) {\textbf{Augur}\\diagnose + predict};
  \node[stp, fill=green!5, right=2em of diag] (branch)
    {\textbf{DAG}\\create BranchNode};
  \node[stp, fill=orange!5, right=2em of branch] (verify)
    {\textbf{F$\star$X}\\verify safety};
  \node[stp, fill=purple!5, below=2em of branch] (eval)
    {\textbf{Evaluate}\\A/B on validation};
  \node[stp, fill=red!5, left=2em of eval] (discard)
    {\textbf{Discard}\\branch removed};
  \node[stp, fill=green!8, right=2em of eval] (activate)
    {\textbf{Keeper}\\atomic swap};

  %% Arrows
  \draw[arr] (diag) -- (branch);
  \draw[arr] (branch) -- (verify);
  \draw[arr] (verify) -- (eval);
  \draw[arr] (eval) -- node[below, lbl] {quality $\uparrow$} (activate);
  \draw[arr] (eval) -- node[below, lbl] {quality $\downarrow$} (discard);

  %% Feedback
  \draw[arr, dashed, black!20] (activate.south)
    .. controls +(0,-1.2) and +(0,-1.2) .. (diag.south)
    node[midway, below=2pt, lbl] {next iteration};
\end{tikzpicture}
\caption{Optimization activation protocol.
  Every optimization follows the same path:
  Augur diagnoses, a BranchNode encodes the change,
  F$\star$X verifies safety, evaluation compares arms on validation data,
  and the Keeper activates or discards via atomic swap.
  No optimization is irreversible.}
\label{fig:optloop}
\end{figure}

Uniform protocol for all optimizations:

\begin{enumerate}
  \item Augur diagnoses and predicts expected improvement.
  \item Optimization implemented as BranchNode or DAG modification.
  \item F$\star$X verifies safety (memory bounds, kernel access).
  \item New branch evaluated on validation data alongside current branch.
  \item Quality maintained + improvement confirmed $\to$ Keeper activates via atomic swap.
  \item Quality degrades $\to$ branch discarded.
\end{enumerate}

Built-in A/B testing with instant rollback.
No optimization is irreversible.
